
\medskip
\phantomsection\label{cor:sga3-VIB-5.6.2}%
\noindent\textbf{Corollary 5.6.2.---}
\emph{Let \(S\), \(G\), \(X\) and \(s\) be as in the theorem, let \(U\) be an open subscheme of \(X\) such
that \(U_s\) is dense in \(X_s\), and let \(A\) be a finite subset of \(X_s\).\footnote{Translator's note: the
1970 text says ``a subset''. The proof uses finitely many points, and the statement fails for infinite subsets:
for \(G=X=\mathbb G_{a,k}\) over a field \(k\), acting by translations, and \(U=X-\{0\}\), no translate
\(b\cdot U_{S''}\) contains the inverse image of the whole fibre. ``Finite'' is the correction of the Gille--Polo
re-edition (2011), Expos\'e~VI\textsubscript{B}. The proof also uses hypothesis~(ii) of the theorem (see the note
there): the points of \(X_s\) specialize to closed points, whose residue fields are finite over \(\kappa(s)\).}
Then there exist a morphism \(g:S''\to S'\to S\) as in Lemma \hyperref[lem:sga3-VIB-5.6.1]{5.6.1} and an element
\(b\in G(S'')\) such that \(g^{-1}(A)\) is contained in \(b\cdot U\).}

After shrinking \(S\), we may suppose that there exists an open neighbourhood \(V\) of the identity section of
\(G\) which is of finite presentation and universally open over \(S\). We shall construct \(b\) in \(V(S'')\). One
sees then that we may again replace \(S\) by the spectrum of a strict henselization of \(\mathcal O_{S,s}\). We
may suppose the points \(a_i\in A\) closed in \(X_s\). If some of the points \(a_i\) have residue fields that are
nontrivial extensions (necessarily finite and radicial) of \(\kappa(s)\), then, after a finite sequence of flat,
finite extensions of finite presentation of the type \(\mathcal O_{S,s}\to\mathcal O_{S,s}[T]/(T^p-x)\)
\((x\in\mathcal O_{S,s})\), we see that we may suppose the \(a_i\) rational. The hypothesis on \(X\) implies that,
for \(i=0,1,\ldots,n\), there exist closed points \(b'_i\) of \(G_s\times X_s\) which are mapped, by the canonical
morphism, onto the points \((a_i,a_0)\) of \(X_s\times X_s\). Proceeding as above, we may suppose the \(b'_i\)
rational, hence of the form \((b_i,a_0)\),\footnote{Translator's note: the 1970 text prints \((b_i,h_0)\).} where
\(b_i\in G_s(\kappa(s))\) is such that \(a_i=b_i\cdot a_0\). Let then \(W_s\) be the open subset of \(G_s\) that
is the inverse image of \(U_s\) under the morphism \(b\mapsto b\cdot a_0\). The hypothesis on \(U_s\) implies that
\(W_s\) is dense in \(G_s\). If \(b_s\) is a rational point of \(G_s\), we then have the equivalences
\[
  a_i\in b_s\cdot U_s\iff b_s^{-1}b_i\cdot a_0\in U_s\iff b_s^{-1}b_i\in W_s\iff b_s\in b_iW_s^{-1}.
\]
Since \(W_s\) is dense in \(G_s\), so are the \(b_iW_s^{-1}\) and their intersection. We can therefore find a
closed point \(x\) of \(V_s\cap\bigl[\bigcap_ib_iW_s^{-1}\bigr]\). After a finite surjective extension of finite
presentation, we may suppose, by Lemma \hyperref[lem:sga3-VIB-5.6.1]{5.6.1},\footnote{Translator's note: the
1970 text calls 5.6.1 ``Lemma 1'' here and in the statement, and at the end of the proof of 5.3 below it calls
5.6.3 ``Lemma 2'' and 5.6.2 ``the corollary to Lemma 1'', numbers of the 1965 edition; the Gille--Polo
re-edition (2011) gives the numbers of this Expos\'e.} that there exists a section \(b\) of \(G\) passing through
\(x\). It is clear that the section \(b\) answers the question.

\par\medskip\noindent
\emph{Addition} (after the Gille--Polo re-edition (2011), Expos\'e~VI\textsubscript{B}; our formulation).
The re-edition proves 5.6.2 by means of a lemma of O.~Gabber, under the sole hypothesis that \(U_s\) is dense in
every irreducible component of \(X_s\), without hypothesis~(ii) of 5.3. Consequently 5.3 holds when (i) and (ii)
are replaced by: (i\('\)) \emph{for every \(s\in S\), there is an open subscheme \(U\) of \(X\), separated over
\(S\), such that \(U_s\) is dense in every irreducible component of \(X_s\)}; this follows from (i) when \(X_s\)
has locally finitely many irreducible components, in particular under (ii). Here is the lemma; \(k\) is a field.
(a) \emph{Let \(G\) be a \(k\)-scheme of finite type, \(K\) an extension of \(k\), \(W\) a dense open subset of
\(G_K\), and \(\pi\colon G_K\to G\) the projection. The points \(g\in G\) such that \(W\) contains every maximal
point of \(\pi^{-1}(g)\) form a dense open subset \(U\) of \(G\).} If \(L\supset K\), the projection
\(G_L\to G_K\), being flat and surjective, maps the maximal points of the fibres of \(G_L\to G\) onto those of
the fibres of \(\pi\); so \(U\) does not change if \(K\) and \(W\) are replaced by \(L\) and the inverse image of
\(W\), and we may suppose that \(K\) contains a separable closure \(k_s\) of \(k\). Let \(p\colon G_K\to G_{k_s}\)
and \(q\colon G_{k_s}\to G\) be the projections. The fibres of \(p\) are irreducible (EGA~IV\(_2\), 4.3.3) and
\(p\) is open, so \(V=p(W)\) is a dense open subset of \(G_{k_s}\), and \(W\) contains the generic point of
\(p^{-1}(g')\) if and only if \(g'\in V\). The fibre \(q^{-1}(g)\) is a finite set; the maximal points of
\(\pi^{-1}(g)\) are the generic points of the \(p^{-1}(g')\), \(g'\in q^{-1}(g)\). Hence \(U\) is the complement
of \(q(G_{k_s}-V)\), which is closed, \(q\) being integral; and \(U\) is dense: \(G_{k_s}-V\), a nowhere dense
closed set, is defined over a finite extension of \(k\), so it has finitely many conjugates under
\(\operatorname{Gal}(k_s/k)\), which acts transitively on the fibres of \(q\), and the complement of their union is
a dense open subset of \(G_{k_s}\) mapped by \(q\) into \(U\).
(b) \emph{Let \(G\) be a geometrically irreducible \(k\)-scheme of finite type, \(\mu\colon G\times_kX\to Y\) a
morphism of \(k\)-schemes, \(x\in X\), and \(\Omega\) an open subset of \(Y\) such that the morphism
\(\mu_x\colon G_{\kappa(x)}\to Y\) induced by \(\mu\) maps some point into \(\Omega\). The points \(g\in G\) such
that \(\mu\) maps every maximal point of the fibre \(g\times x\) of \(G\times_kX\) over \((g,x)\) into \(\Omega\)
form a dense open subset of \(G\); and if \(g\) is a closed point of this set, \(\mu(g\times x)\subset\Omega\).}
Apply (a) with \(K=\kappa(x)\) and \(W=\mu_x^{-1}(\Omega)\), a nonempty, hence dense, open subset of the
irreducible scheme \(G_{\kappa(x)}\); for \(g\) closed, \(\kappa(g)\) is finite over \(k\), so every point of
\(g\times x\) is maximal in it.
(c) \emph{Let \(H\) be a \(k\)-group locally of finite type acting on \(X\) with
\(H\times_kX\to X\times_kX\), \((h,x)\mapsto(hx,x)\), surjective, \(G=H^0\), \(U\) an open subset of \(X\) dense in
every irreducible component of \(X\), and \(A\) a finite subset of \(X\). There is a closed point \(c\) of \(G\)
such that \(c^{-1}a'\in U_{\kappa(c)}\) for every point \(a'\) of \(X_{\kappa(c)}\) above a point of \(A\).}
Indeed \(G\cdot U=X\) by the re-edition's VI\(_A\), \SGARefLink{prop:sga3-VIA-2.6.4}{2.6.4} (in this edition (d)
of the second Addition after VI\(_A\), 2.4.1): every point \(a\) of \(A\) can be written \(a=\gamma u\),
with \(\gamma\) and \(u\) points of \(G\) and \(U\) with values in a field, and then \(\gamma^{-1}a\in U\). So (b),
applied to \(\mu(\gamma,x)=\gamma^{-1}x\), gives for each \(a\in A\) a dense open subset of the geometrically
irreducible scheme \(G\) of finite type (VI\(_A\), 2.4); their intersection contains a closed point \(c\), and
\(c\) answers the question.
To prove 5.6.2, apply (c) to \(G_s\) acting on \(X_s\): it gives a closed point \(c\) of \(G^0_s\) such that
\(c^{-1}\) maps the points above \(A\) into \(U\). As \(G\) is universally open over \(S\) at the generic point of
\(G^0_s\), which lies in every neighbourhood of the identity section, Lemma 5.6.1, applied to \(G\) at \(c\), gives
\(g\colon S''\to S'\to S\) and \(b\in G(S'')\) with \(b(s'')\) above \(c\); then \(b^{-1}a'\in U_{S''}\) for every
\(a'\in g^{-1}(A)\), that is, \(g^{-1}(A)\subset b\cdot U\).
\par\medskip
